%%% Please fill in basic information on your thesis, which will be automatically
%%% inserted at the right places.

% Comment out to exclude papers from the thesis.
\def\includePapers{true}

% Type of your thesis:
%   "bc" for Bachelor's
%   "mgr" for Master's
%   "phd" for PhD
%   "rig" for rigorosum
\def\ThesisType{phd}

% Language of your study programme:
%   "cs" for Czech
%   "en" for English
\def\StudyLanguage{en}

% Thesis title in English (exactly as in the official assignment)
\def\ThesisTitle{A Unified Framework for Multi-Model Data Management and Adaptation}

% Author of the thesis (you)
\def\ThesisAuthor{Mgr. Jáchym Bártík}

% Year when the thesis is submitted
\def\YearSubmitted{2026}

% Name of the department or institute, where the work was officially assigned
% (according to the Organizational Structure of MFF UK in English,
% see https://www.mff.cuni.cz/en/faculty/organizational-structure,
% or a full name of a department outside MFF)
\def\Department{Department of Software Engineering}

% Is it a Department (Katedra), or an Institute (Ústav)?
\def\DeptType{Department}

% Thesis supervisor: name, surname and titles
\def\Supervisor{doc. RNDr. Irena Holubová, Ph.D.}

% Supervisor's department (again according to Organizational structure of MFF)
\def\SupervisorsDepartment{Department of Software Engineering}

\def\AdvisorA{Ing. Pavel Koupil, Ph.D.}

\def\AdvisorADepartment{Department of Software Engineering}

\def\AdvisorB{prof. RNDr. Roman Barták, Ph.D.}

\def\AdvisorBDepartment{Department of Theoretical Computer Science and Mathematical Logic}

\def\MffAddress{%
Faculty of Mathematics and Physics\\
Charles University\\
Malostranské náměstí 25, 11800 Prague 1, Czechia
}

% Study programme (does not apply to rigorosum theses)
\def\StudyProgramme{Computer Science}
\def\StudyBranch{Software Systems}

% Abstract (recommended length around 80-200 words; this is not a copy of your thesis assignment!)
\def\Abstract{%
Multi-model databases combine relational, document, graph, and other data models, but their management remains fragmented across incompatible schemas, query languages, and adaptation mechanisms. Building on an existing category-theoretic foundation for multi-model schemas, mappings, transformations, and querying, this thesis develops methods for propagating conceputal schema changes to logical schemas, and queries, as well as a self-adaptation approach that searches for better logical configurations under changing workloads. The adaptation process combines workload-driven search with learned cost estimation, allowing candidate configurations to be compared only from query plans. The thesis contributes formal models, propagation and synchronization algorithms, and an experimental validation. Together, the results show that multi-model data management can be described coherently, automated to a substantial degree, and adapted to evolving application needs.
}


% 3 to 5 keywords (recommended) separated by \sep
% Keywords are useful for indexing and searching for the theses by topic.
\def\ThesisKeywords{%
multi-model databases\sep query processing\sep schema evolution\sep self-adaptation\sep category theory
}

% If any of your metadata strings contains TeX macros, you need to provide
% a plain-text version for use in XMP metadata embedded in the output PDF file.
% If you are not sure, check the generated thesis.xmpdata file.
\def\ThesisAuthorXMP{\ThesisAuthor}
\def\ThesisTitleXMP{\ThesisTitle}
\def\ThesisKeywordsXMP{multi-model databases, query processing, schema evolution, self-adaptation, category theory}
\def\AbstractXMP{\Abstract}

% If your abstracts are long and do not fit in the infopage, you can make the
% fonts a bit smaller by this setting. (Also, you should try to compress your abstract more.)
\def\InfoPageFont{}
%\def\InfoPageFont{\small}  % uncomment to decrease font size

% If you are studing in a Czech programme, you also need to provide metadata in Czech:
% (in English programmes, this is not used anywhere)

\def\ThesisTitleCS{Unifikovaný framework pro správu a adaptaci multi-modelových dat}
% \def\DepartmentCS{Název katedry česky}
% \def\DeptTypeCS{Katedra}
% \def\SupervisorsDepartmentCS{katedra vedoucího}
% \def\StudyProgrammeCS{studijní program}

\def\ThesisKeywordsCS{%
multi-modelové databáze\sep zpracování dotazů\sep evoluce schématu\sep samo-adaptace\sep teorie kategorií
}

\def\AbstractCS{%
Multi-modelové databáze kombinují relační, dokumentové, grafové a další datové modely. Jejich správa však zůstává roztříštěná napříč nekompatibilními schématy, dotazovacími jazyky a adaptačními mechanismy. Tato práce navazuje na existující framework, založený na teorii kategorií, který obsahuje multi-modelová schémata, mapování, transformace a dotazování. Přitom rozvíjí metody pro propagaci změn konceptuálního schématu do logických shémat a dotazů. Dále navrhuje přístup samo-adaptace, který hledá vhodnější logické uspořádání dat při měnící se zátěži dotazů. Adaptační proces kombinuje vyhledávání řízené zátěží dotazů s naučeným odhadováním ceny dotazů, což umožňuje porovnávat možné konfigurace pouze na základě plánů dotazů. Práce přináší formální modely, propagační a synchronizační algoritmy, a experimentální validaci. Tyto výsledky ukazují, že správu multi-modelových dat lze popsat koherentně, do značné míry automatizovat a přizpůsobovat vyvíjejícím se potřebám aplikací.
}

\def\ThesisTitleCSXMP{\ThesisTitleCS}
